\textbf{5.\ (Pirates and treasure).}
50 pirates ordered by rank are dividing a big treasure that consists of 1000
equal indivisible items. They employ the following procedure. Firstly, the
captain offers a division. If at least half of the crew (including the captain)
agree, then the division is applied. Otherwise the captain is thrown overboard
and the second-ranked pirate offers a division, and so on. Formalize the
situation as a dynamic game and find a subgame perfect Nash equilibrium such
that no pirate is indifferent at any time.

\textbf{Solution.}

Consider a subgame with \(n\) surviving pirates, relabeled
\[
1,2,\dots,n,
\]
where pirate \(1\) is the proposer. A proposal passes if it gets at least
\[
\left\lceil \frac n2 \right\rceil
\]
votes. Since the proposer votes for himself, he needs
\[
\left\lceil \frac n2 \right\rceil -1
=
\left\lfloor \frac{n-1}{2}\right\rfloor
\]
additional votes.

Using backward induction. A pirate votes yes only if he gets strictly more than
he would get after rejecting the proposal. We also don't want to pay the next possible captain and we pay only to those pirates
who would get 0 if a proposal is rejected.

% For small \(n\), the pattern is:
% \[
% n=1:\quad (1000),
% \]
% \[
% n=2:\quad (1000,0),
% \]
% \[
% n=3:\quad (999,1,0),
% \]
% \[
% n=4:\quad (999,1,0,0),
% \]
% \[
% n=5:\quad (998,1,1,0,0)
% \]

For small \(n\), the pattern is:
\[
n=1:\quad (1000),
\]
\[
n=2:\quad (1000,0),
\]
\[
n=3:\quad (999,0,1),
\]
\[
n=4:\quad (999,0,1,0),
\]
\[
n=5:\quad (998,0,1,0,1)
\]


Thus, in a subgame with \(n\) pirates, the proposer buys the cheapest votes:
pirates
\[
3,5,7,\dots
\]
each get \(1\) item, and the others get \(0\).

For \(n=50\), the captain needs
\[
\left\lceil \frac{50}{2}\right\rceil =25
\]
votes total. His own vote counts, so he needs \(24\) more votes.

He gives \(1\) item to pirates
\[
3,5,7,\dots,49.
\]

There are \(24\) such pirates. Therefore the captain keeps
\[
1000-24=976.
\]

Hence the equilibrium proposal is
\[
(x_1,x_2,\dots,x_{50})
=
(976,0,1,0,1,0,\dots,1,0)
\]
where
\[
x_1=976,
\]
\[
x_{3}=x_{5}=\cdots=x_{49}=1,
\]
and
\[
x_2=x_4=\cdots=x_{50}=0
\]

The yes-voters are
\[
\{1,3,5,\dots,49\},
\]
which gives exactly \(25\) votes, so the proposal passes.

Therefore, in the subgame perfect Nash equilibrium, the first captain survives
and keeps
\[
976
\]
items.
